\section{Conclusion}

We introduced \textbf{RiboPO}, a preference-optimization framework for RNA inverse folding that formalizes the problem as preference learning under multiple noisy physical proxies. RiboPO's three load-bearing design choices --- (i) $\varepsilon$-Pareto-dominance preferences on standardized per-metric features, (ii) a structural quality gate that blocks compositional reward hacking, and (iii) frozen-reference DPO with a decreasing-margin curriculum --- map cleanly onto a heteroscedastic Bradley--Terry derivation (Appendix~\ref{app:hetero_derivation}), a trust-region drift bound (Theorem~\ref{thm:trust_region}), and an off-policy bias bound (Theorem~\ref{thm:off_policy}). Ablations show each is individually load-bearing: removing the Pareto/conjunctive filter causes structural collapse; removing the gate causes compositional reward hacking on weaker base models; removing the frozen reference causes unbounded drift. Empirically, RiboPO improves the Pareto front of the SSTT panel: $20/21$ metrics improve, $17$ at $p<0.05$, with $69\%$ of MFE gain GC-independent, $P(\text{target})$ up $687\%$, and a $64\times$ sample-efficiency gain at the joint-criterion pass@k.

\paragraph{Limitations.}
We identify limitations honestly. \emph{(i) Single-class structure predictors:} RhoFold+ during training and the cross-tool panel (NUPACK, RNAstructure, AlphaFold3, DRFold2) at evaluation are not fully independent --- the Turner thermodynamic parameters are shared across NUPACK / Vienna / RNAstructure, and AF3 / DRFold2 / RhoFold+ overlap on PDB training data. Cross-tool consistency is necessary but not sufficient evidence of physical correctness. \emph{(ii) No experimental validation:} all results are computational. SHAPE-MaP / DMS-MaP and UV-melting on a small panel of designed sequences would be the highest-leverage next experiment. \emph{(iii) Ensemble dynamics ignored:} the reward is computed on a single 3D prediction; co-transcriptional folding, kinetic traps, ion-dependent (Mg$^{2+}$) folding, and conformational ensemble interconversion are invisible to RiboPO. \emph{(iv) Static preference pool:} reuse across rounds becomes off-policy (Theorem~\ref{thm:off_policy}); performance plateaus at round 2--4 and gradually degrades for $R \geq 5$. Importance-corrected iterative DPO would extend the round budget. \emph{(v) Recovery trade-off:} the $-5.3\%$ recovery decrease is uniform across structural motifs; INF for non-canonical contacts \emph{improves} ($+24\%$, $p = 1\!\times\!10^{-3}$), suggesting functional preservation, but DSSR base-pair annotation on \emph{predicted} structures is itself noisy.

\paragraph{Future work.}
Several directions follow naturally from the framework: \emph{(a) Multi-objective DPO with explicit Pareto-front parameterization}, learning a weight-conditioned policy $\pi_\theta(s\mid\mathcal{G}, w)$ that traverses the achievable trade-off front at inference time; \emph{(b) Importance-corrected iterative DPO} resampling candidates from $\pi_{\theta, r}$ at each round with truncated importance weights, lifting the round budget beyond the static-pair bound (Theorem~\ref{thm:off_policy}); \emph{(c) Multi-predictor consensus rewards} aggregating $\{\text{RhoFold+}, \text{AF3}, \text{NuFold}\}$ for 3D and $\{\text{Vienna}, \text{NUPACK}, \text{RNAstructure}\}$ for thermodynamics, directly addressing predictor-specific bias; \emph{(d) Conformational-ensemble preferences} computed not from a single prediction but from a structural ensemble (multiple recycles or sampled suboptimal structures), accounting for RNA conformational dynamics; and \emph{(e) Function-conditioned design} that masks recovery loss at residues annotated as functionally critical (riboswitch ligand-binding sites, ribozyme catalytic cores), separating designed-and-functional from designed-and-bland.

The broader contribution of RiboPO is to identify the failure modes that arise when preference optimization is naively transplanted from LLM alignment to a domain with heterogeneous noisy physical proxies and combinatorial sequence-structure degeneracy --- and to provide an analyzable, empirically validated framework that addresses each. We expect the Pareto-dominance preference construction and the heteroscedastic margin to be useful in adjacent biomolecular design problems where multiple imperfect physical proxies must be combined into a single training signal.
